% ============================================================
% Bao cao BTL CO4031 - Phan cua Nguoi B (Dashboard & DSS)
% Bien dich:  xelatex part_B_dss.tex   (chay 2 lan de co muc luc)
% ============================================================
\documentclass[12pt,a4paper]{article}

\usepackage{fontspec}
\setmainfont{Latin Modern Roman}
\setmonofont{DejaVu Sans Mono}[Scale=0.82]
\usepackage[a4paper,margin=2.5cm]{geometry}
\usepackage{setspace}
\onehalfspacing
\usepackage{amsmath}
\usepackage{booktabs,longtable,tabularx,array}
\usepackage{graphicx}
\usepackage{xcolor}
\usepackage{listings}
\usepackage{enumitem}
\usepackage{caption}
\usepackage{tikz}
\usetikzlibrary{arrows.meta,positioning,shapes.geometric}
\usepackage[colorlinks=true,linkcolor=blue!50!black,urlcolor=blue!50!black]{hyperref}

% --- Tieu de tieng Viet ---
\renewcommand{\contentsname}{Mục lục}
\renewcommand{\figurename}{Hình}
\renewcommand{\tablename}{Bảng}
\renewcommand{\listfigurename}{Danh mục hình}
\renewcommand{\listtablename}{Danh mục bảng}
\renewcommand{\refname}{Tài liệu tham khảo}

% --- Code listing ---
\definecolor{codebg}{RGB}{245,245,245}
\lstdefinestyle{sql}{
  language=SQL, basicstyle=\ttfamily\footnotesize, backgroundcolor=\color{codebg},
  keywordstyle=\color{blue!60!black}\bfseries, commentstyle=\color{green!40!black},
  breaklines=true, frame=single, rulecolor=\color{black!20}, showstringspaces=false,
  morekeywords={CASE,WHEN,THEN,ELSE,END,LIMIT,PRAGMA,AUTOINCREMENT,REFERENCES,IF,EXISTS}
}
\lstdefinestyle{py}{
  language=Python, basicstyle=\ttfamily\footnotesize, backgroundcolor=\color{codebg},
  keywordstyle=\color{blue!60!black}\bfseries, commentstyle=\color{green!40!black},
  breaklines=true, frame=single, rulecolor=\color{black!20}, showstringspaces=false
}
\newcommand{\code}[1]{\texttt{\small #1}}
\captionsetup{font=small,labelfont=bf}

\title{\textbf{BÁO CÁO BÀI TẬP LỚN CO4031}\\[4pt]
\Large Kho dữ liệu thương mại điện tử Olist và hệ hỗ trợ ra quyết định\\[4pt]
\large Phần B: Dashboard phân tích và Hệ hỗ trợ ra quyết định (DSS)}
\author{Vỹ (Người B -- Ứng dụng \& DSS)}
\date{Học kỳ 261 -- Năm học 2026--2027}

\begin{document}
\maketitle
\tableofcontents
\clearpage

% ============================================================
\section{Kiến trúc dashboard và DSS}
% ============================================================

\subsection{Vị trí trong hệ thống}
Phần B dùng lại kho dữ liệu (DW) do phần A xây dựng, không đọc lại các tệp CSV gốc. Theo hợp đồng dữ liệu
giữa hai phần (\code{docs/data\_contract.md}), phần B \emph{chỉ đọc} từ tệp \code{data/processed/olist\_dw.db}
đã được ETL của phần A nạp sẵn theo lược đồ star schema (\code{src/sql/01\_schema.sql}). Phần B không được
tự sửa lược đồ hay đọc CSV trực tiếp, để tránh hai phần xung đột dữ liệu.

\begin{figure}[h]
\centering
\begin{tikzpicture}[
  node distance=9mm,
  box/.style={draw, rounded corners, minimum height=1.1cm, minimum width=2.6cm, align=center, fill=blue!6},
  arr/.style={-{Latex[length=2.5mm]}, thick}]
  \node[box, fill=orange!15] (dw) {Kho dữ liệu\\SQLite (star schema)};
  \node[box, below=14mm of dw] (qry) {\code{src/app/queries.py}\\SQL tham số hóa};
  \node[box, below left=14mm and 4mm of qry, fill=green!10] (app) {\code{src/app/app.py}\\Streamlit (5 trang)};
  \node[box, below right=14mm and 4mm of qry, fill=green!10] (dss) {\code{src/dss/}\\forecast, rfm, whatif};
  \draw[arr] (dw) -- (qry);
  \draw[arr] (qry) -- (app);
  \draw[arr] (qry) -- (dss);
  \draw[arr] (dss.north) to[out=100,in=-80] (app.south);
\end{tikzpicture}
\caption{Kiến trúc dashboard và DSS: giao diện và các mô hình DSS đều chỉ truy cập DW qua một tầng truy vấn SQL duy nhất}\label{fig:arch-b}
\end{figure}

\subsection{Các tầng}
\begin{itemize}[nosep]
  \item \textbf{Tầng truy vấn (\code{src/app/queries.py}).} Toàn bộ câu SQL đọc DW nằm trong một tệp duy nhất.
        Mọi điều kiện lọc (năm, bang, danh mục\ldots) dùng tham số \code{?} của \code{sqlite3}, không nối chuỗi,
        tránh SQL injection và giữ quy ước doanh thu \code{SUM(price)}, loại đơn \code{canceled}/\code{unavailable}
        thống nhất ở một chỗ (hàm \code{\_not\_canceled\_clause}).
  \item \textbf{Tầng giao diện (\code{src/app/app.py}).} Ứng dụng Streamlit, chỉ gọi hàm ở \code{queries.py}
        và ba mô-đun DSS, không bao giờ mở kết nối CSDL trực tiếp. Mỗi hàm truy vấn được bọc bằng
        \code{@st.cache\_data} để không truy vấn lại DW mỗi lần người dùng tương tác với widget.
  \item \textbf{Tầng DSS (\code{src/dss/forecast.py}, \code{rfm.py}, \code{whatif.py}).} Ba mô-đun độc lập, có thể
        chạy trực tiếp bằng dòng lệnh (\code{python src/dss/forecast.py}) để kiểm tra kết quả ngoài giao diện,
        đồng thời được \code{app.py} import và gọi lại trong trang \emph{Dự báo \& DSS}.
\end{itemize}

% ============================================================
\section{Dashboard Streamlit}
% ============================================================

Dashboard gồm 5 trang, chọn qua \code{st.sidebar.radio} (hàm \code{main()} trong \code{app.py}, biến \code{PAGES}).

\subsection{Trang 1 -- Tổng quan (\code{page\_overview})}
Hiển thị 5 chỉ số KPI (doanh thu, số đơn, giá trị đơn trung bình, điểm đánh giá trung bình, \% giao trễ) từ
\code{kpi\_summary()}, và biểu đồ đường doanh thu theo tháng từ \code{monthly\_revenue()}.

\begin{figure}[h]
\centering
\fbox{\parbox{0.9\textwidth}{\centering\vspace{3cm}\textit{TODO: chèn ảnh chụp màn hình trang Tổng quan}\vspace{3cm}}}
\caption{Trang Tổng quan của dashboard}\label{fig:ss-overview}
\end{figure}

\subsection{Trang 2 -- Phân tích đa chiều (OLAP) (\code{page\_olap})}
Đây là trang chứng minh năm phép toán OLAP theo yêu cầu đề bài, mỗi phép toán ứng với một hàm riêng trong
\code{queries.py} (Bảng~\ref{tab:olap-map}).

\begin{table}[h]
\centering
\caption{Ánh xạ phép toán OLAP sang hàm truy vấn và thành phần giao diện}\label{tab:olap-map}
\small
\begin{tabularx}{\textwidth}{@{}l l X@{}}
\toprule
\textbf{Phép toán} & \textbf{Hàm (\code{queries.py})} & \textbf{Thao tác trên giao diện}\\
\midrule
Roll-up & \code{rollup\_by\_year()} & Tổng hợp doanh thu từ mức tháng lên mức năm, hiển thị biểu đồ cột theo năm.\\
Drill-down & \code{drilldown\_by\_month(year)} & Chọn một năm (\code{st.selectbox}), xem chi tiết doanh thu theo quý và tháng của năm đó.\\
Slice & \code{slice\_by\_category(year)} & Cố định chiều thời gian ở một năm được chọn, xem doanh thu theo danh mục sản phẩm.\\
Dice & \code{dice\_states\_categories(states, categories, year)} & Chọn nhiều bang và nhiều danh mục (\code{st.multiselect}) cùng lúc, có thể chọn thêm năm, trả về khối con.\\
Pivot & \code{pivot\_category\_by\_year()} & Xoay chiều năm thành cột (\code{y2016}, \code{y2017}, \code{y2018}, \code{total}) bằng biểu thức \code{SUM(CASE WHEN\ldots)}, hiển thị dạng bảng.\\
\bottomrule
\end{tabularx}
\end{table}

\begin{figure}[h]
\centering
\fbox{\parbox{0.9\textwidth}{\centering\vspace{3cm}\textit{TODO: chèn ảnh chụp màn hình trang OLAP (roll-up/drill-down/slice/dice/pivot)}\vspace{3cm}}}
\caption{Trang Phân tích đa chiều (OLAP)}\label{fig:ss-olap}
\end{figure}

\subsection{Trang 3 -- Địa lý \& khách hàng (\code{page\_geo\_customer})}
Biểu đồ cột doanh thu theo bang, tô màu theo điểm đánh giá trung bình (\code{revenue\_review\_by\_state()}),
và biểu đồ tròn cơ cấu doanh thu theo hình thức thanh toán (\code{payment\_mix()}).

\begin{figure}[h]
\centering
\fbox{\parbox{0.9\textwidth}{\centering\vspace{3cm}\textit{TODO: chèn ảnh chụp màn hình trang Địa lý \& khách hàng}\vspace{3cm}}}
\caption{Trang Địa lý \& khách hàng}\label{fig:ss-geo}
\end{figure}

\subsection{Trang 4 -- Giao hàng \& hài lòng (\code{page\_delivery})}
So sánh điểm đánh giá trung bình và số ngày giao trung bình giữa nhóm giao trễ và đúng hạn
(\code{delivery\_vs\_review()}), rồi tách theo 10 bang và 10 danh mục có số lượng dòng hàng lớn nhất
(\code{delivery\_vs\_review\_by\_state()}, \code{delivery\_vs\_review\_by\_category()}).

\begin{figure}[h]
\centering
\fbox{\parbox{0.9\textwidth}{\centering\vspace{3cm}\textit{TODO: chèn ảnh chụp màn hình trang Giao hàng \& hài lòng}\vspace{3cm}}}
\caption{Trang Giao hàng \& hài lòng}\label{fig:ss-delivery}
\end{figure}

\subsection{Trang 5 -- Dự báo \& DSS (\code{page\_dss})}
Gồm ba khối tương ứng ba mô hình DSS trình bày ở Mục~3: biểu đồ dự báo doanh thu kèm bảng MAE/MAPE, bảng
phân khúc khách hàng RFM + K-Means kèm biểu đồ elbow, và hai khối what-if (giao trễ, khuyến mãi) với thanh
trượt tương tác (\code{st.slider}).

\begin{figure}[h]
\centering
\fbox{\parbox{0.9\textwidth}{\centering\vspace{3cm}\textit{TODO: chèn ảnh chụp màn hình trang Dự báo \& DSS}\vspace{3cm}}}
\caption{Trang Dự báo \& DSS}\label{fig:ss-dss}
\end{figure}

% ============================================================
\section{Mô hình DSS}
% ============================================================

Ba mô-đun trong \code{src/dss/} đều tuân thủ hợp đồng dữ liệu: chỉ nhận dữ liệu đầu vào từ
\code{src/app/queries.py}, không tự mở kết nối CSDL khác. Số liệu dưới đây lấy từ kết quả chạy thực tế
(\code{python src/dss/forecast.py}, \code{rfm.py}, \code{whatif.py}) trên kho dữ liệu đã nạp.

\subsection{Dự báo doanh thu (\code{src/dss/forecast.py})}

\paragraph{Dữ liệu vào.} Chuỗi doanh thu theo tháng, 20 điểm từ \code{2017-01} đến \code{2018-08}
(\code{monthly\_revenue\_series\_for\_forecast()}). Khoảng thời gian này được cắt có chủ đích: năm 2016 chỉ
có 3 tháng rời rạc và tháng 9/2018 chỉ có 1 đơn, đưa vào sẽ làm nhiễu mô hình.

\paragraph{Phương pháp.} Giữ lại 4 tháng cuối (\code{2018-05}..\code{2018-08}) làm tập kiểm tra, 16 tháng đầu
làm tập huấn luyện (\code{TEST\_SIZE = 4}). So sánh hai mô hình:
\begin{itemize}[nosep]
  \item \textbf{Baseline (trung bình trượt):} dự báo phẳng, bằng trung bình 3 tháng cuối của tập huấn luyện,
        lặp lại cho toàn bộ chân trời dự báo (\code{\_naive\_moving\_average\_forecast}, \code{window=3}).
  \item \textbf{Holt (xu hướng tuyến tính):} \code{statsmodels.tsa.holtwinters.ExponentialSmoothing} với
        \code{trend="add"}, \code{seasonal=None} -- \emph{không} có thành phần mùa vụ vì chuỗi quá ngắn để
        ước lượng đáng tin cậy. Mô hình được huấn luyện lại trên toàn bộ 20 tháng để cho ra dự báo 1 bước kế
        tiếp, kèm khoảng tin cậy 95\% tính từ độ lệch chuẩn phần dư (\code{resid\_std}).
\end{itemize}

\paragraph{Kết quả thực tế.} (Bảng~\ref{tab:forecast-result})

\begin{table}[h]
\centering
\caption{Sai số dự báo trên tập kiểm tra (4 tháng 2018-05..2018-08)}\label{tab:forecast-result}
\small
\begin{tabular}{@{}l r r@{}}
\toprule
\textbf{Mô hình} & \textbf{MAE} & \textbf{MAPE (\%)}\\
\midrule
Baseline (trung bình trượt) & 69.432,04 & 7,85\\
Holt (xu hướng tuyến tính) & 281.179,04 & 32,24\\
\bottomrule
\end{tabular}
\end{table}

Dự báo tháng kế tiếp sau \code{2018-08} (mô hình Holt, huấn luyện lại trên toàn bộ chuỗi): \textbf{854.929,55},
khoảng tin cậy 95\%: \textbf{643.623,52 .. 1.066.235,57}.

\paragraph{Giải thích vì sao Holt tệ hơn baseline.} Đúng như cảnh báo trong docstring của mô-đun: doanh thu
tăng gần như liên tục rồi đạt đỉnh tại tháng 11/2017 (Black Friday, khoảng 1,0 triệu, xem phần A), sau đó
\emph{chững lại và giảm nhẹ}. Mô hình Holt ước lượng một xu hướng tăng tuyến tính từ 16 tháng huấn luyện (kết
thúc ở \code{2018-04}, thời điểm doanh thu vẫn đang hồi phục sau Tết), nên ngoại suy vượt quá doanh thu thực
tế của 4 tháng kiểm tra tiếp theo -- trong khi baseline (trung bình trượt phẳng) tình cờ bám sát mức doanh thu
tương đối ổn định của giai đoạn này hơn. Đây không phải lỗi cài đặt mà là hệ quả trực tiếp của chuỗi thời gian
ngắn, chưa đủ 2 chu kỳ mùa vụ để mô hình xu hướng được kiểm chứng đáng tin cậy.

\subsection{Phân khúc khách hàng RFM + K-Means (\code{src/dss/rfm.py})}

\paragraph{Cái bẫy F (Frequency) và cách xử lý.} Trong 95.420 khách hàng (theo \code{customer\_key} đã gộp
theo người thật), chỉ khoảng 2.913 người (\textasciitilde 3\%) có từ 2 đơn trở lên. Vì vậy Frequency theo
cách tính kinh điển gần như luôn bằng 1 và không thể chia phân vị (\code{qcut}) một cách có ý nghĩa. Giải
pháp của mô-đun: thay F bằng hai đặc trưng khác --
\textbf{\code{repeat\_flag}} (cờ nhị phân: 1 nếu khách có \(\ge\) 2 đơn, ngược lại 0) và
\textbf{\code{avg\_review}} (điểm đánh giá trung bình của khách), vì điểm đánh giá mang tín hiệu hành vi thật
mà một cột tần suất luôn bằng 1 không thể cung cấp.

\paragraph{Dữ liệu vào và đặc trưng.} Từ \code{customer\_order\_history()} (1 dòng / đơn hàng), gộp theo
\code{customer\_key} thành 4 đặc trưng: \code{recency\_days} (số ngày từ lần mua gần nhất đến ngày tham chiếu
\textbf{2018-09-03}, một tháng sau đơn cuối trong cửa sổ phân tích), \code{monetary} (tổng chi tiêu),
\code{repeat\_flag}, \code{avg\_review} (NULL được điền bằng trung vị của khách khác). Dữ liệu được chuẩn hóa
bằng \code{StandardScaler} trước khi phân cụm.

\paragraph{Chọn K.} Chạy K-Means với \(K \in [2,6]\), 10 lần khởi tạo mỗi K (\code{n\_init=10}). Vì
\code{silhouette\_score} có độ phức tạp \(O(n^2)\), tính trên mẫu con 5.000 khách hàng. Kết quả (Bảng~\ref{tab:elbow}):
K=2 có silhouette cao nhất (0,681) nhưng mô-đun giới hạn lựa chọn trong khoảng 3--4 theo yêu cầu đề bài
(\code{chosen\_k = max(3, min(4, chosen\_k))}) để có đủ số phân khúc cho hành động nghiệp vụ; trong khoảng đó
K=3 có silhouette tốt nhất.

\begin{table}[h]
\centering
\caption{Inertia và silhouette theo K}\label{tab:elbow}
\small
\begin{tabular}{@{}r r r@{}}
\toprule
\textbf{K} & \textbf{Inertia} & \textbf{Silhouette}\\
\midrule
2 & 283.974,6 & 0,681\\
3 & 208.709,2 & 0,411\\
4 & 153.541,4 & 0,399\\
5 & 108.683,3 & 0,422\\
6 & 95.277,0 & 0,430\\
\bottomrule
\end{tabular}
\end{table}

\textbf{K đã chọn: 3, silhouette = 0,411.}

\paragraph{Hồ sơ phân khúc.} (Bảng~\ref{tab:segments}) Việc đặt tên phân khúc theo thứ tự ưu tiên trong
\code{\_name\_segments}: cụm có tỷ lệ mua lại cao nhất (và chi tiêu cao nhất) được gọi là \emph{Khách giá trị
cao}; trong các cụm còn lại (chủ yếu khách mua 1 lần, nơi recency/monetary gần như không khác biệt -- chính
\code{avg\_review} là yếu tố phân tách), cụm có điểm đánh giá cao nhất là \emph{Khách mới} (hài lòng ở đơn đầu,
đáng để chăm sóc cho đơn thứ hai), cụm có điểm đánh giá thấp nhất là \emph{Khách ngủ đông}.

\begin{table}[h]
\centering
\caption{Hồ sơ từng phân khúc khách hàng}\label{tab:segments}
\small
\begin{tabular}{@{}l r r r r r r@{}}
\toprule
\textbf{Phân khúc} & \textbf{Recency (ngày)} & \textbf{Monetary} & \textbf{Repeat flag} & \textbf{Review TB} & \textbf{Số khách} & \textbf{\%}\\
\midrule
Khách ngủ đông & 248,5 & 149,5 & 0,000 & 1,89 & 20.300 & 21,4\\
Khách mới & 242,5 & 134,8 & 0,000 & 4,75 & 71.792 & 75,6\\
Khách giá trị cao & 226,6 & 271,1 & 0,999 & 4,17 & 2.891 & 3,0\\
\bottomrule
\end{tabular}
\end{table}

Nhóm \emph{Khách giá trị cao} (3,0\% khách hàng) gần như trùng khít với nhóm khách mua lại (\code{repeat\_flag}
trung bình 0,999) và có chi tiêu cao nhất. Hai nhóm còn lại đều là khách mua 1 lần (repeat\_flag = 0) nhưng
tách biệt rõ theo mức độ hài lòng: \emph{Khách mới} hài lòng (4,75/5, chiếm 75,6\% khách hàng) trong khi
\emph{Khách ngủ đông} không hài lòng (1,89/5, 21,4\%).

\paragraph{Hành động đề xuất.} Khách giá trị cao: chăm sóc VIP, ưu đãi độc quyền, giao hàng ưu tiên. Khách
mới: gửi voucher cho đơn thứ 2, gợi ý sản phẩm liên quan để biến thành khách mua lại. Khách ngủ đông: khảo sát
lý do không hài lòng, chiến dịch win-back/giảm giá để tránh mất khách vĩnh viễn.

\subsection{Mô hình what-if (\code{src/dss/whatif.py})}

Mô-đun cung cấp hai mô hình có \textbf{tình trạng tri thức (epistemic status) khác nhau}, phải được phân biệt
rõ trên giao diện và trong báo cáo.

\paragraph{What-if 1: Giảm giao trễ \(\to\) điểm hài lòng -- dựa trên số đo thực tế.}
Dựa trên quan hệ \emph{đã đo được} trong DW (phần A, Bảng~giao trễ): đơn giao trễ có điểm đánh giá trung bình
\textbf{2,55}, đơn đúng hạn/sớm \textbf{4,21}. Mô hình giả định hai điểm trung bình này không đổi, chỉ tỷ
trọng hai nhóm thay đổi, rồi tính trung bình có trọng số
(\code{weighted\_avg = p*2.55 + (1-p)*4.21}). Đây là phép nội suy tuyến tính đơn giản, \emph{không phải} mô
hình nhân quả đầy đủ. Tỷ lệ giao trễ hiện tại trên dữ liệu: \textbf{7,91\%}. Ví dụ kết quả (Bảng~\ref{tab:whatif-delivery}).

\begin{table}[h]
\centering
\caption{What-if giảm tỷ lệ giao trễ (hiện tại: 7,91\%)}\label{tab:whatif-delivery}
\small
\begin{tabular}{@{}r r r@{}}
\toprule
\textbf{Tỷ lệ trễ mục tiêu (\%)} & \textbf{Điểm TB dự kiến} & \textbf{\(\Delta\) so với hiện tại}\\
\midrule
7,91 (hiện tại) & 4,079 & +0,000\\
4,91 & 4,128 & +0,050\\
1,91 & 4,178 & +0,100\\
0,00 & 4,210 & +0,131\\
\bottomrule
\end{tabular}
\end{table}

\paragraph{What-if 2: Khuyến mãi theo danh mục \(\to\) doanh thu -- dựa trên giả định người dùng.}
Người dùng nhập mức giảm giá (\code{discount\_pct}) và \textbf{độ co giãn cầu theo giá}
(\code{elasticity}, quy ước dương = cầu tăng khi giá giảm). Công thức:
\(\%\Delta\text{quantity} = \text{elasticity} \times \%\Delta(-\text{price})\), từ đó
\(\text{Revenue}_{new} = \text{Revenue}_{base} \times (1-d) \times (1 + \text{elasticity} \times d)\) với
\(d\) là mức giảm giá. \textbf{Độ co giãn là giả định do người dùng tự đặt, không phải số đo thực tế hay dự
báo chắc chắn}; giao diện luôn hiển thị cảnh báo này. Ví dụ với danh mục \code{health\_beauty}
(doanh thu nền: 1.255.695,00, Bảng~\ref{tab:whatif-promo}):

\begin{table}[h]
\centering
\caption{What-if khuyến mãi -- danh mục health\_beauty}\label{tab:whatif-promo}
\small
\begin{tabular}{@{}r r r r@{}}
\toprule
\textbf{Giảm giá (\%)} & \textbf{Co giãn} & \textbf{Doanh thu dự kiến} & \textbf{\(\Delta\) (\%)}\\
\midrule
10 & 1,0 & 1.243.138 & $-1,00$\\
10 & 2,0 & 1.356.151 & $+8,00$\\
20 & 1,5 & 1.305.923 & $+4,00$\\
\bottomrule
\end{tabular}
\end{table}

Với co giãn = 1,0 (cầu tăng đúng bằng \% giảm giá), doanh thu vẫn giảm nhẹ vì phần tăng số lượng không bù đủ
phần giảm giá trên từng đơn vị; cần co giãn \(>\) 1 (cầu nhạy hơn giá) thì khuyến mãi mới làm tăng doanh thu
trong mô hình này.

% ============================================================
\section{Đánh giá}
% ============================================================

\begin{table}[h]
\centering
\caption{Tổng hợp chỉ số đánh giá các mô hình DSS}\label{tab:eval-summary}
\small
\begin{tabularx}{\textwidth}{@{}l l X@{}}
\toprule
\textbf{Mô hình} & \textbf{Chỉ số} & \textbf{Giá trị}\\
\midrule
Dự báo -- Baseline & MAE / MAPE & 69.432,04 / 7,85\%\\
Dự báo -- Holt & MAE / MAPE & 281.179,04 / 32,24\% (kém hơn baseline)\\
RFM -- K-Means & K đã chọn / Silhouette & 3 / 0,411\\
What-if giao trễ & Căn cứ & Số đo thực tế (2,55 vs 4,21)\\
What-if khuyến mãi & Căn cứ & Giả định người dùng (độ co giãn)\\
\bottomrule
\end{tabularx}
\end{table}

\paragraph{Hạn chế của từng mô hình (nêu trung thực, không phóng đại kết quả).}
\begin{itemize}[nosep]
  \item \textbf{Dự báo doanh thu.} Chuỗi chỉ có 20 tháng, chưa đủ 2 chu kỳ mùa vụ (24 tháng trở lên) nên không
        thể ước lượng và kiểm chứng một mô hình mùa vụ đáng tin cậy; vì vậy nhóm không dùng Holt-Winters đầy
        đủ mà chỉ dùng Holt không mùa vụ. Trên dữ liệu hiện tại, Holt có sai số \emph{cao hơn} baseline vì
        xu hướng tăng ước lượng từ giai đoạn huấn luyện không còn đúng ở giai đoạn kiểm tra (doanh thu chững
        lại sau đỉnh Black Friday 11/2017). Khoảng tin cậy 95\% của dự báo tháng kế tiếp khá rộng
        (643.623,52 -- 1.066.235,57), phản ánh đúng mức độ bất định với chuỗi ngắn.
  \item \textbf{RFM + K-Means.} Silhouette 0,411 ở K=3 là mức \emph{trung bình}, không phải phân cụm tách biệt
        rõ ràng; K=2 có silhouette cao hơn (0,681) nhưng bị giới hạn lại trong khoảng 3--4 theo yêu cầu đề bài
        để có đủ phân khúc hành động. F (Frequency) đã được thay bằng cờ nhị phân mua lại + điểm đánh giá
        trung bình vì \textasciitilde 97\% khách hàng chỉ có 1 đơn, nên đây không phải RFM kinh điển mà là một
        biến thể phù hợp với đặc điểm dữ liệu.
  \item \textbf{What-if giao trễ.} Giả định điểm trung bình của mỗi nhóm (trễ/đúng hạn) không đổi khi tỷ trọng
        thay đổi -- là một phép nội suy, không phải mô hình nhân quả (không loại trừ các yếu tố gây nhiễu khác
        ảnh hưởng đến điểm đánh giá).
  \item \textbf{What-if khuyến mãi.} Độ co giãn cầu theo giá là tham số \emph{người dùng tự nhập}, chưa được
        ước lượng từ dữ liệu lịch sử (DW không có đủ các đợt khuyến mãi có ghi nhận giá gốc/giá giảm để ước
        lượng co giãn thực tế). Kết quả chỉ mang tính minh họa kịch bản, không phải dự báo.
\end{itemize}

% ============================================================
\section{Hạn chế và hướng phát triển}
% ============================================================

\subsection{Hạn chế hiện tại}
\begin{itemize}[nosep]
  \item \textbf{Chưa có ảnh chụp màn hình dashboard.} Báo cáo này dùng khung \code{\textbackslash fbox} làm chỗ
        trống (TODO) cho 5 trang dashboard; cần chụp và chèn ảnh thật trước khi nộp/thuyết trình.
  \item \textbf{Chuỗi thời gian ngắn.} 20 tháng dữ liệu không đủ để kiểm định mô hình mùa vụ hoặc các mô hình
        dự báo phức tạp hơn (ARIMA mùa vụ, Prophet\ldots).
  \item \textbf{Mô hình Holt hiện đang kém hơn baseline} trên tập kiểm tra 4 tháng cuối -- là một kết quả thật,
        không phải lỗi cài đặt, nhưng cho thấy cần thận trọng khi dùng Holt để ra quyết định thực tế ở giai
        đoạn dữ liệu hiện tại.
  \item \textbf{F (Frequency) trong RFM bị thay thế} bằng cờ nhị phân mua lại vì đặc thù dữ liệu (hầu hết khách
        chỉ mua 1 lần trong giai đoạn quan sát), nên các phân khúc không hoàn toàn tương đương RFM kinh điển.
  \item \textbf{Độ co giãn cầu trong what-if khuyến mãi} là tham số giả định, chưa được kiểm chứng bằng dữ liệu
        thực nghiệm (A/B test hay lịch sử khuyến mãi).
\end{itemize}

\subsection{Hướng phát triển}
\begin{enumerate}[nosep]
  \item Thu thập/tích lũy thêm dữ liệu theo thời gian để chuỗi doanh thu đạt tối thiểu 24 tháng, cho phép dùng
        Holt-Winters có thành phần mùa vụ và so sánh với các mô hình như SARIMA hoặc Prophet.
  \item Ghi nhận các đợt khuyến mãi thực tế (giá gốc, giá giảm, doanh số trước/sau) để ước lượng độ co giãn cầu
        theo giá từ dữ liệu thay vì để người dùng tự giả định.
  \item Bổ sung ảnh chụp màn hình và video demo thực tế cho từng trang dashboard; cân nhắc thêm bản đồ theo
        mã bưu chính nếu phần A mở rộng dùng dữ liệu địa lý chi tiết.
\end{enumerate}

\end{document}
